% Appendix C, part 3 of 3: the model, transport, universality, local Mellin inversion. Last part of appendix C.

\subsection{The model}
\label{sec:evi-model}

\Cref{tab:evi-model} lists the main checks of \cref{sec:mod}.

\begin{table}[ht]
\centering\small
\begin{tabular}{>{\raggedright\arraybackslash}p{0.4\textwidth}|>{\raggedright\arraybackslash}p{0.36\textwidth}|>{\raggedright\arraybackslash}p{0.16\textwidth}}
statement & range & arithmetic\\ \hline
topological recursion on $x=\Wv$ is $(-5)^{h-1}$ times the model CohFT, i.e.\ the model correlators are $(-5)^{1-h}$ times those of topological recursion (\cref{thm:mod-tr}) & $72120$ colour-basis correlators with $\psi$-deficit $\leq 6$ at $(h,n)=(0,3)$, $(0,4)$, $(0,5)$, $(1,1)$, $(1,2)$, $(1,3)$, $(2,1)$, $(2,2)$, $(3,1)$ ($4665$ non-zero); on $x=\kappa\Wv$, $\kappa=1,4,-1$: $6495$ each & two primes\\
the R-matrix of the curve $x=\kappa\Wv$ is $\Mo(z/(\kappa\varsigma^4))$ (\cref{prop:mod-Rtr}) & all $25$ entries through $z^{10}$, $\kappa=1,4,-1$, $\varsigma=1,2$ & exact, $\bbQ(\zeta)$\\
the dilaton-leaf sign of \cite[Def.~2.13]{PPZ15} & $(h,n)=(1,1)$, four points, $\kappa=1,4$; the opposite sign fails at all four & exact\\
$\FTR_h(v=1)=B_{2h}/(2h(2h-2))=-\tfrac{1}{240},\tfrac{1}{1008},-\tfrac{1}{1440},\tfrac{1}{1056}$ & $h=2,\dots,5$ & $\bbQ(\zeta)$ ($h\leq 3$); two primes\\
the period identity of \cref{prop:mod-period}; $\Sigma_h=-4^hB_{2h}/(2h)$ & $h\leq 7$; $h\leq 20$ & exact\\
the lattice theorem (\cref{thm:mod-lattice}) & weight $\leq 400$; $\leq 1000$; $16200$ raw-leg kernel vectors & exact; two primes\\
the leg dictionary (\cref{prop:mod-LD}) & $35000$ pairings per prime; $52$ frame one-point functions, $h=2,3$ & two primes; $\bbQ(\zeta)$\\
the Model Lemma on $\sK$ (\cref{thm:mod-MLK}) & $5289$ correlators with $h\leq 4$, including $1302$ at $(h,n)=(2,4)$; $1395$ singular, all zero & exact; two primes\\
\end{tabular}
\caption{Checks on the model.}
\label{tab:evi-model}
\end{table}

Further, modulo two primes unless stated: the limits $v\to 0$ of the correlators at $13$ pairs $(h,n)$, for instance $\omega_{1,1}\to-\tfrac58\dd p/p^6$ and $\omega_{2,1}\to\tfrac{2079}{128}\dd p/p^{16}$; the closed formulas of the correlators against the recursion, exactly at $(0,3)$, $(1,1)$, $(0,4)$, $(1,2)$, $(2,1)$; the formula of \cref{rmk:mod-hock} for $h\leq 12$ (exactly). The $28$ non-zero pairings of legs of the class $\cL$ at $(1,1)$, $(2,1)$, $(1,2)$, $(2,2)$ have $v$-exponent $\geq 0$, while legs outside $\cL$ are singular (for example the pairing of $\omega_{1,1}$ with $-\log p$ is $-\tfrac1{12}v^{-1}$). The residues $\Res_{\varsigma_b}x^k\omega_{1,1}$ vanish for $b=0,\dots,4$ and $k=0,\dots,6$, while $40$ of the $45$ residues of $p^k\omega_{1,1}$, $k=1,\dots,9$, do not, and $\Res_{p=0}x^5\omega_{1,1}=-15/4$, so that the circle pairing of $x^5$ is not zero (\cref{def:mod-pair}). In the base case of \cref{thm:mod-V} the local formula of \cref{lem:mod-w11} gives $-1/60$ at each branch point at $v=1$, in total $-1/12$ (the opposite orientation of the ray gives $+1/12$), and the integral of the closed form of $\omega_{1,1}$ gives $-1/(12v)$ at $v=1,32,-3$. No log-free leg lies outside $\cL$ among $10400$ (weight $\leq 320$) and $16200$ (weight $\leq 400$) vectors, exactly; \cref{prop:mod-dii} holds for $k\leq 100$ and $n\leq 2000$ exactly and $n\leq 3000$ modulo two primes; \cref{prop:mod-V1} holds on $32559$ tuples with $h\leq 4$ ($21733$ by a second implementation), and the wrong-sign controls fail. The singular lists of the Model Lemma at genus $0,1,2,3$ (at most $5,5,4,3$ insertions, frames with $k\leq 3h+3$, $\psi$-deficit $\leq 6$; $43$, $418$, $425$, $127$ lists) all vanish: $326$ by genuine cancellation, $642$ with every term zero although allowed, $45$ trivially; deliberately wrong frames give non-zero values. The $17416$ singular ordered tuples of legs in $\cL$ and $-\log p$ at $(4,2),(3,3),(2,4),(1,5),(0,6)$ vanish modulo each of the two primes, and so do $40739$ multisets with $h+n\leq 7$, $h\leq 6$, by a second implementation.

\emph{Off $\sK$} (\cref{rmk:mod-offK}), exactly at $v=1$: $\langle\psi^5H\rangle_3=53/1209600=-\tfrac{53}{192}(\Fm_3)'$, whereas the one-point law of \cref{cor:mod-onepoint} would give $\tfrac1{96}(\Fm_3)'=-1/604800$; at genus $4$ the remainders of $\psi^5H$, $\psi^6\,1$, $\psi^7H^4$ are $\tfrac{19}{100800000}(1,-1,-\tfrac12)$; and the quotient $(\langle XY\rangle_h-(\Fm_h)''\Pi_0(X)\Pi_0(Y))/(\Fm_h)'$ depends on $h$, for instance $-23/12$ at $h=2$ and $-13/48$ at $h=3$ for $X=\psi H$, $Y=\psi^2H^3$.

\subsection{Transport and the comparison}
\label{sec:evi-transport}

\Cref{lem:tra-E} holds exactly over $\bbQ$ for random symplectic $R=\exp(X)$, generic and $\bbZ/5$-graded, with random $f$ in both slots: the numerator is divisible by $\psi+\psi'$ if and only if (Sp) holds, and the two sides agree; two further implementations, with $R=\exp(X)$ and with $R$ the Cayley transform of a random $X$, modulo $u^8$ and $u^7$, for $f$ built from $\YC$ and for random $f$ with the required decay, agree on all coefficients. Controls fail: a sign flip, a non-symplectic $R$, and an $f$ without decay, which gives different values for different truncations. For $\Pc$, in both frames, $z^0,\dots,z^{16}$, $u<16$, exactly: (Sp); (G) on $575$ complete coefficients; $T_P=\tc$ through $u^{14}$; $\Pi_0(q_T)=\tc$ and $\Pi_1(q_T)=\dots=\Pi_7(q_T)=0$; the edges divisible by $\psi+\psi'$ and annihilated by $\Pi_0,\dots,\Pi_7$ on $2985$ coefficients (\cref{cor:tra-edges}, \cref{lem:tra-T}); in the other normalisation (G) fails on $77$ coefficients, $\Pi_5(q_T)\neq 0$, the edge condition fails on $92$, and symplecticity fails from $z^{10}$. For the outer rows $R'$ of the natural form ($\zh\leq 24$, $u<18$, both primes, both frames) the edges are annihilated by $\Pi_k$ for $k=0,\dots,20$ ($2160$ coefficients each) and the leaf point by $\Pi_1,\dots,\Pi_{20}$ modulo $u^{18}$; with the scalar $\cn_5=1/7$ every condition fails and the leaf point acquires exactly a $\Pi_5$ component. In both frames $\delta_P=dT_P-v=u^2/30+(127/3150)u^3+\dots$ for $P=\Pc$ (it does not depend on the frame: $T_{\Pc}=\tc$ by \cref{prop:tra-can} and $d=1/p_{0,0}=\tau^{1/4}\varpi_0$ by \cref{lem:gau-diag}, so $dT_{\Pc}=\tau^{1/4}\varpi_1$); hence $\tc-\hat v=\delta_P/d=u^2/30+O(u^3)$, $\hat v=v/d$, the value used in \cref{rmk:red-leaf}.

\emph{Frames.} $\YC$ from the closed form of \cref{lem:gau-vanshape} agrees with the regular solution of the model system ($n<96$); the frames with $k\leq 80$, built with $\Theta^{-1}=\sum_r(-N)^r$ (\cref{lem:tra-theta}) and by elimination, agree ($324$ frames), satisfy $\Pi_i(\fr(k,a))=0$ for $0\leq i\leq k+12$, and have integral exponents $\geq 1$; by generic elimination at $v=1,2$ for $k\leq 60$ the trivial frames are exactly the ten with $k<\nu_a+1$ (\cref{lem:tra-frames}), and the rule $k\leq\nu_a+1$ is refuted. Controls fail: frames built from $\YC(+z)$ ($304$ of $314$ not annihilated) and $\YC$ without its terms $z^{-n}$, $n\geq 5$ ($214$ of $224$). In the $u$-adic setting, modulo $u^{10}$ (\cref{lem:tra-conv}): $\Theta$ is invertible, frame decompositions reconstruct random elements, $3312$ orders of frame coefficients equal $(4(k-j)+\nu_a+1)/5$, and a test bilinear form converges with the sharp bound $4K/5-3$; with $v=1+u$ in place of $v(u)$ there is no $u$-adic limit.

\emph{Theorem~R and the comparison.} In the flat sum of $\Pc$ (\cref{def:red-flatsum}) at genus $2$, frame $(0,0)$, the vertex graph is $u^{-2}(\tfrac1{1200}-\tfrac{7u}{6000}+\tfrac{287u^2}{864000}+\dots)$ and the six edge graphs are regular, with the values of \cref{tab:evi-g2}; in the frame $(1/3,-2/7)$ the same holds with the polar part $(1/1200)\tc^{-2}$. Controls: removing the leaves gives $-\tfrac1{900}u^{-1}$; the dressings $(1+u^mzE_{21})P$, $m=0,1$, are singular from order $u^m$; the non-gauge dressings $\exp(u^mzN)P$, $N=E_{04}+E_{43}$, are regular, with $T_P=u+u^2/5+\dots$ for $m=0$ and $T_P=\tc+O(u^3)$ for $m=1$. The comparison theorem (\cref{thm:flh-main}) at genus $3$: the flat and the outer terms of the $42$ graphs agree on the coefficients of $u^{-4},\dots,u^3$ ($336$ per prime and frame; frame $(0,0)$ modulo both primes, frame $(1/3,-2/7)$ modulo $2^{89}-1$), of which $108$ coefficients of edge graphs and $4$ of the vertex graph are non-zero; a second implementation reproduced all $252$ non-zero coefficients of the regrouping. Controls change many coefficients: a constant $\Mo$ in the splitting ($94$ of $252$), the dilaton factor dropped ($77$ of $336$), the leaves removed ($88$), the sign of the edge dropped ($103$), the other normalisation ($21$, in six graphs), $R(-z)$ for $R(z)$ ($112$), the dressing $\exp(uzN)$ ($74$).

\subsection{Universality and necessity}
\label{sec:evi-univ}

\emph{The other normalisation} (\cref{rmk:mai-stored}, \cref{rmk:mai-graphlocal}). Its scalar series $\cn(z)=P_{11}(0,z)=1+m_5z^5+m_{10}z^{10}+\dots$ has, exactly over $\bbQ$ and in both frames, $m_5=-43397/807975000$ and $m_{10}=-765368951400317/1305647201250000000$, so that $\cn(z)\cn(-z)=1+(2m_{10}-m_5^2)z^{10}+\dots$ with $2m_{10}-m_5^2=-18223115111903/15543419062500000$. At genus $4$ exactly six edge graphs are singular, the two one-edge bridges (between vertices of genus $1$ and $3$, and of genus $2$ and $2$) from order $u^4$ of $G_\Gamma$, and at order $u^5$ the genus-$3$ vertex with a loop and three graphs with two edges, while the other $372$ terms are unchanged and the polar part of the total does not move. Per unit $m_5$ the responses at the orders $u^4$, $u^5$ are $-\tfrac{1}{8400000},-\tfrac{187}{680400000}$ for the vertex graph, $-\tfrac{1}{1008000},\tfrac{13}{6804000}$ and $\tfrac{1}{900000},-\tfrac{53}{32400000}$ for the two bridges, and $-\tfrac{41}{68040000},-\tfrac{11}{34020000},\tfrac{49}{16200000},-\tfrac{17}{8100000}$ for the four graphs at $u^5$; they sum to zero at both orders, and the responses at $u^4$ are the same for perturbed systems. Two implementations agree.

\emph{Universality} (\cref{cor:mai-univ}, \cref{rmk:mai-annih}, \cref{lem:mai-coord}). Canonical gauges of perturbed one-pole systems that do not come from a Picard--Fuchs operator make all $378$ edge graphs of genus $4$ regular, and the vertex graph is then $(1/180000)(T/\lambda)^{-6}$ plus a regular series; likewise six perturbed systems at genus $2$ and $3$, in the system's own coordinate $T$. Annihilating non-gauge dressings are regular at genus $2$ ($6$ graphs), $3$ ($41$) and $4$ ($378$), with the vertex graph $c_gT_{\mathrm{eff}}^{2-2g}$, for instance $T_{\mathrm{eff}}=u+\tfrac{49}{45}u^2+\dots$ against $\tc=u+\tfrac{7}{10}u^2+\dots$; dressings that violate annihilation or flat regularity, and random orbits of the first column, give singular edge graphs (for example $14$ at genus $3$). Changes of the product $\Bprod$: for eight systems all edge graphs are regular at genus $2$ and $3$ modulo both primes, and at genus $4$ (those with at most two edges, and all $378$ for two systems), with the vertex graph in the system's own coordinate; $\phi'^5\Bprod_X(\phi)=\Bprod_q$ holds exactly for the change of coordinate $\phi$ of \cref{lem:mai-coord}; the controls fail as predicted. At first order, at genus $2$ ($9$ variables) the dressings that keep every edge graph regular are exactly the annihilating ones; at genus $3$ ($30$ variables) annihilation implies regularity, and the two further regular directions break the vertex law and are singular at genus $4$; without flat regularity ($20$, resp.\ $68$, variables) annihilation no longer implies regularity ($5$, resp.\ $23$, independent responses survive).

\subsection{Local Mellin inversion and the formal lemmas}
\label{sec:evi-lmi}

\emph{Theorem LMI on general data} (\cref{thm:gam-lmi}), exactly over $\bbQ$: for $\varphi=0$, $\cF=1$ the series $c^{\cF}$ is $1$ ($M_*=-2,\dots,3$, through $z^{12}$); in ten random cases ($a_*\in\{-1,2,-1/3,5/2,5,-7/4\}$, $M_*\in\{-1,0,1,2\}$) and four cases with entire data, the one-variable series equals the chart series away from $s_*$, the chart at $s_*$ equals \eqref{eq:gam-lmi}, and the Laurent expansion at $s_*$ has no pole (through $z^{12}$ and $z^{14}$); twelve further random cases (among them the non-integer values $M_*=1/3,-2/5,7/2$, outside \cref{def:gam-lmidata}; cf.\ \cref{rmk:gam-Mcomplex}) and a case with $\varphi''(s_*)=0$, through $z^{10}$ exactly and $z^{12}$, $z^{14}$ modulo two primes, and four more cases modulo these primes, two of them with $a_*=\sqrt{-1}$; the two-variable chart equals the one-variable series coefficientwise in $z$ and in $L-s_*$. \Cref{lem:gam-deltasrc} on random data with $c\in\{0,\pm 1,2,-3,1/5\}$ through $z^{12}$, the moment table against the Wick recursion up to degree $24$; the wrong sign of $\delta$, a term $m^3$ and a term $s^2m$ break it. For non-integer $M_*=1/5,1/3,-1/2,7/4$, with the Bernoulli-polynomial form of \cref{rmk:gam-Mcomplex}, the series is regular with the shifted value, through $z^{14}$.

\emph{Falsifiers} (\cref{rmk:gam-necessity}~(c)), exactly, $n\leq 12$: $\Gamma(a/z)^2$ gives poles of order exactly $n$ at $z^n$, with leading coefficients $((2n-1)!!)^2(-1/a_*)^n/(n!\,16^n)$ ($\tfrac1{16},\tfrac9{512},\tfrac{75}{8192},\tfrac{3675}{524288}$ at $a_*=-1$); $\Gamma(a/z)\Gamma(-a/z)$ poles of order $2n$; $\Gamma(2a/z)$ and $\Gamma(a/z)\Gamma(1+a/z)$ poles of order $n$; three coincident factors the leading coefficients $\tfrac19,\tfrac2{81},\tfrac{14}{2187},\tfrac8{19683},-\tfrac{440}{177147}$; a perturbation of the coefficient of $(z/a)^{n_0}$ in the degenerate Stirling series a first pole at $z^{n_0}$, of order $n_0$, with coefficient the perturbation divided by $a_*^{n_0}$; an exponent $a^E$, $E\neq M_*-1/2$, a pole at $z^1$ with coefficient $(1/24+\beta_1(M_*)-E^2/2)/a_*$; $\cF=1/(\sigma-s_*)$ poles of order $n$ with coefficient $n!(-1/a_*)^n$ (observed); $\cF(s_*,0)=0$ with $M_*\neq 0$ simple poles; a degenerate argument with a quadratic term coefficients depending on $\log\mathsf Q$; an amplitude depending on $\lambda$ a logarithmic divergence. Positive controls hold: $\cF$ depending holomorphically on $\mathsf Q$; the exponent $+1/2$ with $M_*=0$ (that is $\Gamma(1+w)=w\Gamma(w)$); $(\sigma-s_*)\Gamma(a/z)$.

\emph{The quintic.} In the form of \cref{thm:gam-lmi} at $s_*=1$: no pole and $c(z;1)=h$, exactly through $z^{12}$ and modulo two primes through $z^{16}$, also at the colours $1,2$ with an explicit $\zeta$ modulo a prime, where the value $h(z)$ in place of $h(\zeta^{-i}z)$ fails. With $\Gamma(5\sigma/z)$ in place of $\Gamma(1+5\sigma/z)$ the value at $L=1$ is still $h$, while the $c_n(L)$ change from $n=2$ on (\cref{rmk:gam-blind}); with $\Gamma(1+(1-\sigma)/z)$, that is the measure $\dd x_0$ for the degenerate factor, the series is regular with the value for $\delta=z$, which differs from $h$ at $z^1$; at $L=0$ with $\Gamma(5\sigma/z)$ the value is $\exp(5\sum_{5\mid 2k-1}\bk_kz^{2k-1})\neq\Ginf$; at $L=0$ the Bernoulli polynomials at $0$ and $-1/5$ in place of $1/5$ fail at $z^5$, and at $4/5$ only at $z^{10}$; the twist weight $+5$ in $h$ fails at $z^1$.

\emph{Other targets} (\cref{rmk:gam-general}), exactly: the $\cO(5)$-twisted $\bbP^4$ at four sets of distinct rational weights, all $20$ colours, through $n=24$, with the values $h_{\bw_i}$, $\bw_i=\{\lambda_i-\lambda_j\}_{j\neq i}\cup\{-5\lambda_i\}$; $\bbP^1$ and $\bbP^2$ (two choices of weights each), local $\bbP^2$ (two lifts), $\cO(-1)\oplus\cO(-1)$ over $\bbP^1$ (two lifts), $\cO(3)$, $\cO(4)$, $\cO(6)$ over $\bbP^4$ and $\cO(4)$ with a shifted lift ($12$ theories, $40$ colours, through $n=20$), with the values predicted from the twisted I-functions, in which an Euler-class twist contributes minus its fibre weight and an inverse Euler class plus its fibre weight; the charts at $q=\infty$; random regular data.

\emph{The seven-variable chart} (\cref{rmk:gam-sevenvar}). The chart series equals $c(z;L)$ exactly over $\bbQ(\zeta)$ through $z^3$ at $L=2,-2/3$ and modulo the primes $2^{62}+177$ and $2^{80}+85$ through $z^6$ at $L=3,5/4,-1/2$; at $L=1$ it equals $h$ (exactly through $z^7$); the seven-variable series equals $c(z;L)$ exactly through $z^4$ at $L=2$; without the Jacobian, or with $\dd y/y$, the chart fails at $z^1$, resp.\ $z^2$. Expanded in $q$, the chart reproduces $h$ times the I-function of colour $0$ on all $98$ non-zero coefficients of $q^dz^n$ with $d\leq 6$, $-d\leq n\leq 10$, exactly.

\emph{The formal lemmas of \cref{sec:fsp}} on random data, exactly: \cref{lem:fsp-ibp}, \cref{lem:fsp-cov}, \cref{lem:fsp-fubini} (inner dimensions $1$ and $2$), the formal delta lemma invoked in \cref{rmk:gam-sevenvar} (which is not stated in \cref{sec:fsp}) and \cref{lem:fsp-param} through $z^6$ to $z^8$; \cref{lem:fsp-stirling} for $m=-3,\dots,3$, both measures, at $a=1,-2/7,3/2$, through $z^{40}$ (through $z^{80}$ at $a=1$); \cref{lem:fsp-dlmf5118} as polynomials in $\theta$ through $w^{-60}$, and $B_k(\theta+1)-B_k(\theta)=k\theta^{k-1}$ for $k\leq 61$; the non-integer form of Euler's integral behind \cref{rmk:gam-Mcomplex}, with a symbolic exponent, through $z^{12}$.
